\def \Subject {مدل‌های گرافی احتمالاتی}
\def \Session {یک}
\renewcommand{\chaptername}{بخش}
\chapter{\Subject}

\section{مقدمه}
در این فصل به تفصیل نظریه و کاربردهای مدل‌های گرافیِ احتمالاتی (PGMs) را بررسی می‌کنیم. هدف ارائهٔ ابزارهای مفهومی و عملی برای توصیف ساختار وابستگی در داده‌ها، روش‌های استنتاج دقیق و تقریبی، و الگوریتم‌های یادگیری پارامتر و ساختار است.

\section{مبانی نظری و استقلال شرطی}
در مدل‌های گرافی، مفهوم استقلال شرطی محوری‌ترین نقش را بازی می‌کند. برای سه متغیر $X,Y,Z$، نوشتهٔ $X \perp\!\!\perp Y\mid Z$ نشان می‌دهد که $X$ و $Y$ با شرط بر $Z$ مستقل‌اند؛ یعنی
$$
p(X,Y\mid Z)=p(X\mid Z)p(Y\mid Z).
$$
این قیدها به ما اجازه می‌دهد تا توزیع‌های پیچیده را به حاصل‌ضرب مؤلفه‌های ساده‌تر تجزیه کنیم.

\subsection{d-separation در شبکه‌های جهت‌دار}
برای یک DAG، قاعدهٔ d-separation تعیین می‌کند که کدام گره‌ها مستقل شرطی هستند. سه ساختار پایه وجود دارد: زنجیره ($A\to B\to C$)، منشعب ($A\leftarrow B\to C$) و هم‌گرا ($A\to B\leftarrow C$). در دو مورد اول، با شرط‌گذاری بر گرهٔ میانی استقلال ایجاد می‌شود، ولی در حالت هم‌گرا شرط‌گذاری روی گرهٔ همگرا یا فرزندان آن وابستگی را ایجاد می‌کند.

\section{عامل‌بندی و نمایش توزیع}
در شبکه‌های بیزی توزیع کلی به صورت ضرب توزیع‌های شرطی محلی نوشته می‌شود:
$$
p(X)=\prod_{i} p(X_i\mid \mathrm{Pa}(X_i)).
$$
در شبکه‌های بدون جهت (MRF) نمایش با عامل‌ها است:
$$
p(X)=\frac{1}{Z}\\prod_{c\in C} \phi_c(X_c),
$$
که محاسبهٔ ثابت نرمال‌سازی $Z$ معمولاً هزینه‌بر است و علت استفاده از روش‌های تقریبی می‌باشد.

\section{استنتاج دقیق: حذف متغیر و پیام‌رسانی}
الگوریتم حذف متغیر (Variable Elimination) برای محاسبهٔ احتمال شرطی $p(X_a\mid X_b)$، به ترتیب متغیرهای غیرمرتبط را حذف کرده و عامل‌های جدیدی تولید می‌کند. پیچیدگی فضایی و زمانی وابسته به treewidth گراف است.

الگوریتم پیام‌رسانی (Belief Propagation) بر روی گراف‌های درختی دقیق است؛ روی گراف‌های چرخه‌دار به صورت تقریبی اجرا می‌شود. پیام‌ها از یک گره $i$ به $j$ به صورت زیر تعریف می‌شوند:
$$
m_{i\to j}(x_j)=\\sum_{x_i} \phi_{i,j}(x_i,x_j)\\prod_{k\in N(i)\\setminus j} m_{k\to i}(x_i).
$$

\section{استنتاج تقریبی}
روش‌های متداول شامل:
\begin{itemize}
    \item نمونه‌برداری MCMC (مثلاً Gibbs sampling): مناسب برای فضای حالت‌های پیوسته یا گسسته اما کند در همگرایی.
    \item تقریب واریانسی: با حداقل‌سازی KL بین یک خانوادهٔ سادهٔ توزیع‌ها و توزیع هدف، که سریع‌تر و مقیاس‌پذیرتر است.
    \item روش‌های مبتنی بر مهمان‌گیری (Importance Sampling): برای برآورد انتگرال‌ها در حالت‌های خاص استفاده می‌شود.
\end{itemize}

\section{یادگیری پارامترها}
با فرض ساختار معلوم، پارامترها (مثلاً جداول احتمال شرطی در شبکه‌های بیزی) را می‌توان با حداکثر درستنمایی (MLE) یا برآورد بیزی محاسبه کرد. در حضور دادهٔ گمشده، الگوریتم EM برای تقریب MLE کاربردی است:
\begin{enumerate}
    \item E-step: محاسبهٔ امید لگ-احتمال کامل نسبت به توزیع پنهان با پارامترهای فعلی.
    \item M-step: بیشینه‌سازی امید حاصل و بروز رسانی پارامترها.
\end{enumerate}

\section{یادگیری ساختار}
یادگیری ساختار شامل جستجو در فضای گراف‌ها است. رویکردها عبارت‌اند از جستجوی حریصانه (greedy search)، الگوریتم‌های مبتنی بر ساختار-پذیرش (score-based)، و روش‌های عامل-محور (constraint-based) مانند PC algorithm که از آزمون‌های استقلال شرطی استفاده می‌کنند. معیارهای نمره‌دهی رایج: BIC، MDL و امتیازات بیزی.

\section{پیاده‌سازی و ملاحظات پیچیدگی}
پیچیدگی محاسباتی وابسته به ساختار گراف است. treewidth کوچک مشکل را قابل حل می‌کند؛ در غیر این صورت روش‌های تقریبی یا تقریب ساختار باید در نظر گرفته شوند. برای کاربردهای عملی توجه به مقیاس‌پذیری، حافظه و زمان اجرا ضروری است.

\section{مثال عددی کوتاه}
فرض کنید سه متغیر باینری $A,B,C$ با ساختار $A\to B\to C$ داریم. اگر $p(A=1)=0.2$, $p(B=1\mid A=1)=0.9$, $p(B=1\mid A=0)=0.1$, $p(C=1\mid B=1)=0.8$, $p(C=1\mid B=0)=0.2$، آنگاه $p(C=1)$ با جمع‌بندی متغیرهای داخلی به‌دست می‌آید:
$$
p(C=1)=\\sum_{a,b} p(a)p(b\mid a)p(C=1\mid b).
$$
محاسبهٔ مستقیم نتیجهٔ عددی را به خواننده می‌سپاریم تا روش جمع‌بندی را تمرین کند.

\section{جمع‌بندی}
این فصل دیدگاه عمیقی نسبت به ساختارها، قواعد استقلال شرطی، روش‌های استنتاج و یادگیری در PGMs ارائه کرد. کاربردها و محدودیت‌های محاسباتی نیز تشریح شد که زمینهٔ مناسبی برای درک پیوند میان این نظریه و روش‌های مولد عمیق فراهم می‌آورد.

\def \Subject {عنوان/سوال/یک}
\def \Session { یک}
% \setcounter{chapter}{\Session-1}
\renewcommand{\chaptername}{سوال}
\chapter{\Subject}

مقدمه سوال اول.

\section{عنوان/بخش}
توضیحات
\begin{itemize}
    \item    
     پوشش کامل مطالب
    \item    
     رعایت قواعد نگارشی دستور زبان فارسی
    \item    
     استفاده از اشکال مناسب
    \item    
     اشاره به منابع کمک آموزشی
\end{itemize}


برای نوشتن شبه‌کد
 \footnote{pseudocode}
 می‌توانید از مثال زیر در الگوریتم
  \ref{pseudocode1}
   استفاده کنید :

\begin{latin}
\begin{algorithm}[H]
    \KwData{this text}
    \KwResult{how to write algorithm with \LaTeX2e }
    initialization\;
    \While{not at end of this document}{
        read current\;
        \eIf{understand}{
            go to next section\;
            current section becomes this one\;
        }{
            go back to the beginning of current section\;
        }
    }
    \label{pseudocode1}
    \caption{How to write algorithms}
\end{algorithm}
\end{latin}

همچنین برای اضافه کردن شکل می‌توانید از شکل زیر استفاده کنید و برای ارجاع دادن به بصورت شکل
 \ref{fig:dynamicprogramming}
 استفاده کنید.
 همچنین برای درج کلمات انگلیسی در پاراگراف فارسی می‌توان به این شکل
 \lr{English Text}
 یا برای تاکید به این شکل
 \grayBox{English Text}
 عمل کرد.
\begin{figure}[H]
    \centering
    \includegraphics[width=0.4\linewidth]{images/dynamic_programming}
    \caption{عنوان شکل}
    \label{fig:dynamicprogramming}
\end{figure}

برای گنجاندن قطعه‌ای از کد به زبان‌های 
\grayBox{Python} 
از فایل اصلی کد می‌توانید به شکل زیر عمل کنید.
 

با مقایسه متن لاتک و خروجی به این مساله پی‌ خواهید برد که برای اشاره به یک قطعه کد در لاتک لازم است از دستور
\grayBox{label}
داخل محیط/بخش 
\grayBox{program}
استفاده کرده و برای اشاره به آن بخش از کد از دستور
\grayBox{programref}
استفاده نمایید. لازم است پارامتر داده شده به هر دو دستور یکی باشد تا بدرستی به قطعه کد مورد نظر اشاره کنید(\programref{pythonfib}). 

\begin{program}[H]
\inputminted[frame=none]{python}{sample.py}
\caption{تابع فیبوناچی در پایتون }
\label{pythonfib}
\end{program}

در نهایت استفاده از دستور 
\grayBox{grayBox}
نیز می‌تواند کمک شایانی به زیبایی و خوانایی متن شما بکند. این دستور علاوه بر عوض کردن رنگ پس‌زمینه از فونت انگلیسی با عرض ثابت نیز استفاده می‌کند که برای کلمات کلیدی یا اسامی متغیرها در کد قابل استفاده است.

\begin{program}[H]
\begin{minted}[frame=none]{python}
class Node: 
def __init__(self, data): 
    self.value = data
    self.next = None  
\end{minted}
\label{pythonnode}
\caption{تعریف لیست پیوندی در ‌پایتون}
\end{program}

\grayLBox{
    AAC\\
    ACG\\
    GAA\\
    GTT\\
    TCG\\
}
\\
\newpage
هم‌چنین نتایج خود را در قالب جدول به فرم زیر گزارش کنید:
\begin{table}[h!]
  \begin{center}
    \caption{عنوان جدول}
    \label{tab:table1}
    \begin{tabular}{l|S|r}
      \textbf{ مدل ۱} & \textbf{مدل ۲} & \textbf{مدل ۳}\\
      $\alpha$ & $\beta$ & $\gamma$ \\
      \hline
      1 & 1110.1 & a\\
      2 & 10.1 & b\\
      3 & 23.113231 & c\\
      4 & 25.113231 & d\\ % <-- added row here
    \end{tabular}
  \end{center}
\end{table}
\\
\\
\\
همچنین سعی کنید حتی‌الامکان به منابع و مراحع مناسب ارجاع دهید
\cite{CLRS}. 
علاوه بر مراجع چنانچه ابزار یا وب‌سایت قابل توجهی موجود است خوب است به آن هم ارجاع دهید
\cite{visualizationwebsite}. 
